% Independently audited fragment; the parent supplies theorem environments.
\section{Checking a template through its two instances}
\label{sec:template-completion}

The transport part of discharge supplies a formula template $F(v)$,
two terms $s,t$, and checked proofs of $s=t$ and $N(F(s))$.
Its proposed result is $N(F(t))$, which must equal the already checked
outer conclusion.  A separate traversal checking the scope and normality
of $F$ is redundant.  This statement concerns only this template: all
actual equations and conclusion formulas remain sorted, scoped, and
normal-form checked.

\begin{lemma}[Completion of an accepted template]
Suppose the generic substitution traversal of a finite JSON expression
$F$ succeeds at both checked numerical terms $s,t$.  Suppose the two
resulting normalized expressions are well-scoped formulas in the local
numerical environment (the second may satisfy the stricter outer scope).
Then there is a well-scoped canonical formula template $F'(v)$, using
only that environment and its one formal parameter $v$, such that
\[
                  N(F'(s))=N(F(s)),\qquad
                  N(F'(t))=N(F(t))
\]
as literal syntax.
\end{lemma}

\begin{proof}
First consider sorting.  At every proper child the generic traversal
uses the exact expected numerical or formula sort, before performing
any erasing reduction.  At its root it may temporarily accept an
unclassified list.  A root of length other than one, two, or three
does not reduce and cannot produce a checked formula.  A numerical
root, including a recursor or an integer replaced by a numerical term,
stays numerical under every reduction.  Thus the successful checked
formula instance forces the original root to be an equation or a paired
formula.  The input is consequently a sorted formula template, possibly
with extra numerical parameters or reducible numerical subterms.

Interpret every negative level outside the allowed local range as an
additional free numerical parameter.  At binder depth $d$, interpret a
positive index larger than $d+1$ in the template as another free numerical
parameter: the one allowed index above the local binders is the formal
parameter $v$.  This interpretation is consistent under traversing a
binder because the excess index is measured relative to its depth.
There are finitely many additional parameters in $F$.

Let $\rho$ replace all of these additional free parameters by $0$, fixing
$v$, the permitted local parameters, and every bound variable.  Neither
$s$ nor $t$ contains an additional parameter.  Capture-avoiding
substitution and normalization coherence therefore give, for $u=s,t$,
\[
 N\bigl((\rho F)[u/v]\bigr)
   =N\bigl(\rho(F[u/v])\bigr)
   =N\bigl(\rho N(F[u/v])\bigr)
   =N(F[u/v]).
\]
The last equality holds because the checked endpoint contains no
additional parameter and is already normal.  In particular, allowing a
new reduction after closing an extra parameter does not invalidate the
argument; coherence, rather than a claim about identical reduction
traces, is being used.

Finally set $F'=N(\rho F)$.  It is sorted, well scoped, and canonical.
A second application of normalization--substitution coherence gives the
two required literal identities.  The one-hole template may contain
zero, one, or several occurrences of $v$; none of the argument requires
linearity.
\end{proof}

Replacing the supplied template by $F'$ leaves the two premises and the
exported conclusion unchanged.  Hence every successfully checked use of
the relaxed interface is an instance of ordinary equality transport.
Conversely, every previously checked canonical template is still accepted.
This removes one trusted traversal without changing the consequence
relation or permitting an out-of-scope active formula.

The equation premise must be checked before the template is substituted,
so that $s,t$ are known to be numerical terms.  The branch instance and
the outer conclusion must both retain their existing checks.  Omitting
one of these conditions is not justified by this lemma.
